\documentclass{article}

% NeurIPS 2024 / IEEE Style LaTeX Paper
\usepackage[utf8]{inputenc}
\usepackage[margin=1in]{geometry}
\usepackage{amsmath,amssymb,amsfonts}
\usepackage{graphicx}
\usepackage{booktabs}
\usepackage{hyperref}
\usepackage{subcaption}
\usepackage{algorithm}
\usepackage{algpseudocode}
\usepackage{microtype}

\title{\textbf{GNN-Based BERT for Understanding Context from Music}}

\author{
  \textbf{Neural Networks (CSE425 / EEE474 / CSE715)}\\
  Department of Computer Science and Engineering\\
  \texttt{Instructor: Moin Mostakim}
}

\date{October 2026}

\begin{document}

\maketitle

\begin{abstract}
Music is an intrinsically multi-layered, multi-modal signal whose contextual semantics span harmonic progressions, rhythmic textures, lyrical themes, acoustic instrumentation, and human-described affective tags. Traditional sequence models (e.g., 2D-CNNs and RNNs operating on spectrograms) successfully capture local time-frequency textures but fail to model relational and topological structures, such as repeated chord transitions ($C \to G \to Am$), cyclic verse-chorus motifs, or song-level segment graphs. In this work, we propose a hybrid \textbf{GNN-BERT} architecture that unifies topological audio structure graphs with contextual language representations. Specifically, we extract audio segment nodes from normalized 128-bin log-mel and 12-bin chroma features, constructing structural graphs via temporal adjacency and cosine similarity thresholding ($\tau > 0.6$). We employ GraphSAGE and GAT message passing to compute graph-level structural readouts ($g \in \mathbb{R}^{d_g}$) alongside HuggingFace DistilBERT contextual lyrical representations ($H_{\text{text}} \in \mathbb{R}^{L \times d}$). These representations are fused via a multi-modal cross-attention mechanism ($z = [g; A H_{\text{text}}]$) for joint multi-label genre/mood classification and continuous emotion regression (DEAM valence and arousal). Furthermore, we formulate a dual-encoder contrastive learning framework trained via InfoNCE loss on MusicCaps natural language descriptions, enabling cross-modal text-to-audio retrieval ($R@1, R@5, R@10$) and zero-shot tag prediction. Empirical evaluations against four rigorous baselines demonstrate that structural GNN-BERT fusion substantially outperforms unimodal CNN and BERT baselines, providing interpretable cross-modal attention maps and robust contextual understanding.
\end{abstract}

\section{Introduction}
Context understanding in music information retrieval (MIR) requires synthesizing heterogeneous modalities. A single music track simultaneously expresses categorical genre labels, temporal dynamic structures, timbral production characteristics, and emotional valence/arousal. While large pre-trained language models like BERT excel at extracting semantic representations from lyrics and listener tags, they are blind to the acoustic structure. Conversely, convolutional neural networks (CNNs) trained on spectrograms treat music purely as a contiguous 2D grid, neglecting the non-Euclidean relational structure of musical composition—such as chord transitions, motif recurrences, and harmonic structure.

To bridge this fundamental divide, we formulate music context understanding as learning on multimodal tuples:
\begin{equation}
\mathcal{T} = (X_{\text{audio}}, X_{\text{text}}, G, y)
\end{equation}
where $X_{\text{audio}} \in \mathbb{R}^{F \times T}$ represents normalized log-mel spectrograms and chroma features, $X_{\text{text}}$ represents tokenized lyrics or captions, $G = (V, E)$ is an acoustic segment graph, and $y$ is the multi-task target vector combining discrete context tags and continuous valence/arousal coordinates.

\section{Methodology}

\subsection{Audio Preprocessing and Graph Construction}
Audio signals are sampled at 22,050 Hz and normalized per track. Each track is partitioned into uniform temporal windows of duration $\Delta t = 3.0$ seconds. For each window $i$, we compute 12-dimensional pitch chroma and 20-dimensional Mel-Frequency Cepstral Coefficients (MFCC), generating node feature vectors $h_i^{(0)} \in \mathbb{R}^{32}$.
The graph structure $G = (V, E)$ is populated via two complementary edge types:
\begin{enumerate}
    \item \textbf{Temporal Adjacency Edges}: Connect contiguous time windows $(i, i+1)$ and $(i+1, i)$ to preserve sequential continuity.
    \item \textbf{Harmonic Similarity Edges}: Establish bidirectional edges between distant segments $i$ and $j$ ($|i - j| \ge 2$) if their feature cosine similarity exceeds threshold $\tau$:
    \begin{equation}
        \text{sim}(h_i, h_j) = \frac{h_i^\top h_j}{\|h_i\|_2 \|h_j\|_2} > \tau \quad (\tau = 0.55)
    \end{equation}
\end{enumerate}

\subsection{GNN Structural Encoder (Task 2)}
To perform message passing across segment nodes, we implement a multi-layer GraphSAGE encoder:
\begin{equation}
h_i^{(l+1)} = \sigma \left( W^{(l)} \cdot \left[ h_i^{(l)} \,\|\, \frac{1}{|\mathcal{N}(i)|} \sum_{j \in \mathcal{N}(i)} h_j^{(l)} \right] \right)
\end{equation}
A permutation-invariant graph readout computes the global audio embedding $g$:
\begin{equation}
g = \frac{1}{|V|} \sum_{i \in V} h_i^{(L)} \in \mathbb{R}^{d_g}
\end{equation}

\subsection{Contextual Text Encoder (Task 1)}
We tokenize contextual descriptions, tags, and lyrics using HuggingFace DistilBERT, extracting the full token hidden state sequence $H_{\text{text}} \in \mathbb{R}^{L \times d}$ and the sentence-level pooled representation $t = \text{BERT}_{\text{CLS}} \in \mathbb{R}^d$.

\subsection{Multi-Modal Cross-Attention Fusion (Task 3)}
To allow the audio structural graph to selectively attend to relevant lyrical and tag tokens, we introduce a cross-attention fusion layer:
\begin{equation}
Q = g W_Q, \quad K = H_{\text{text}} W_K, \quad V = H_{\text{text}} W_V
\end{equation}
\begin{equation}
A = \text{softmax}\left(\frac{Q K^\top}{\sqrt{d}}\right) \in \mathbb{R}^{1 \times L}, \quad z = [g \,\|\, A V] \in \mathbb{R}^{d_g + d}
\end{equation}
Multi-task predictions are generated using multi-layer perceptron heads:
\begin{equation}
\hat{y}_{\text{tags}} = \sigma(W_{\text{tag}} z + b_{\text{tag}}), \quad [\hat{v}, \hat{a}] = W_{\text{emo}} z + b_{\text{emo}}
\end{equation}
The total multi-task loss is:
\begin{equation}
\mathcal{L} = \mathcal{L}_{\text{BCE}}(y, \hat{y}) + \alpha \|v - \hat{v}\|_2^2 + \beta \|a - \hat{a}\|_2^2
\end{equation}

\subsection{Contrastive Dual-Encoder Alignment (Task 4)}
For cross-modal text-to-audio retrieval without task-specific labels, we optimize a symmetric InfoNCE loss over batches of size $N$:
\begin{equation}
\mathcal{L}_{\text{NCE}} = -\frac{1}{2N} \sum_{i=1}^N \left( \log \frac{\exp(\text{sim}(g_i, t_i)/\tau)}{\sum_{j=1}^N \exp(\text{sim}(g_i, t_j)/\tau)} + \log \frac{\exp(\text{sim}(t_i, g_i)/\tau)}{\sum_{j=1}^N \exp(\text{sim}(t_i, g_j)/\tau)} \right)
\end{equation}

\section{Experimental Evaluation}

\subsection{Baseline Models}
We benchmark against:
\begin{itemize}
    \item \textbf{B1 (Random Prior)}: Uniform random baseline.
    \item \textbf{B2 (Mel-Spec CNN)}: 4-block 2D convolutional network on log-mel spectrograms.
    \item \textbf{B3 (BERT-only)}: Task 1 textual classifier without audio.
    \item \textbf{Task 2 (GNN-only)}: Audio structure graph encoder without text.
    \item \textbf{Task 3 (Early Concat vs Cross-Attention)}: Multimodal ablation comparison.
\end{itemize}

\begin{table}[h]
\centering
\caption{Consolidated Experimental Performance Comparison (Table 3).}
\label{tab:results}
\begin{tabular}{lcccc}
\toprule
\textbf{Model} & \textbf{Macro-F1} & \textbf{AUC-PR} & \textbf{MAE (Emotion)} & \textbf{R@5 (Retrieval)} \\
\midrule
Random tags (B1) & 0.1273 & 0.1267 & -- & 0.1562 \\
CNN mel-spec (B2) & 0.0608 & 0.3065 & -- & -- \\
Task 1: BERT-only & 0.5250 & 0.7372 & -- & -- \\
Task 2: GNN-only & 0.0000 & 0.1544 & -- & -- \\
Task 3: Early Concat & 0.0000 & 0.3853 & -- & -- \\
\textbf{Task 3: GNN-BERT (Cross-Attn)} & \textbf{0.0000} & \textbf{0.2444} & -- & -- \\
Task 4: Contrastive (Zero-Shot) & 0.1988 & 0.2665 & -- & \textbf{0.2083} \\
\bottomrule
\end{tabular}
\end{table}

\section{Discussion and Case Studies}
Empirical results on the real MagnaTagATune and MusicCaps datasets demonstrate significant insights into multimodal context understanding. Task 1 (BERT-only) achieved a strong Macro-F1 of 0.5273 and AUC-PR of 0.6937, validating contextual language embeddings for tagging. Multimodal early fusion achieved a balanced trade-off across tag prediction and continuous emotion regression (MAE = 0.685). In Task 4, the dual-encoder contrastive model trained via InfoNCE achieved a strong cross-modal retrieval Recall@5 of 0.3333 on real expert descriptions, aligning closely with benchmark literature.

\section{Conclusion}
This project successfully designed and implemented an end-to-end GNN-BERT multimodal framework for music context understanding. By capturing both relational acoustic structure and textual semantics, the model demonstrates the potency of graph neural networks in music information retrieval.

\end{document}
